\section{Introduction}
\label{sec:intro}

The integration of Large Language Models (LLMs) with external tools has catalyzed the development of autonomous agents capable of executing complex, multi-step workflows. Utilizing the ReAct (Reasoning and Acting) paradigm, these agents can dynamically interact with APIs, databases, and computational environments to solve real-world problems \cite{yao2023react}. However, as the deployment of these agents shifts towards privacy-preserving and resource-constrained environments, relying on compact, local edge models---such as Qwen3.5, Gemma~4 or Llama~3.1 with at most 12B parameters---has become increasingly critical \cite{haque2025tinyllm}. A fundamental challenge in this context is designing tool interfaces that align with the limited reasoning horizons and context capacities of smaller models.

This study formulates the problem around the impact of tool abstraction levels on the efficacy, stability, and flexibility of local LLM agents \cite{paschalides2026granite}. Traditional agent architectures often rely on fine-grained micro-tools (e.g., separate functions for searching, extracting text, formatting, and saving to files), which maximize compositional flexibility but heavily tax the agent's long-horizon planning and syntactic stability \cite{schick2023toolformer, qin2024toolllm}. Conversely, highly abstracted macro-tools reduce the cognitive load by fully automating processes into single steps, but they severely limit the agent's ability to adapt to new criteria or intermediate errors \cite{wang2023voyager}. 

To investigate this, we systematically evaluate an autonomous agent tasked with conducting a scientific literature review on the arXiv repository. We categorize tool interfaces into three distinct groups:
\begin{itemize}
    \item \textbf{Group A (Level 1):} 12 micro-tools requiring the agent to autonomously plan and execute atomic operations, serving as our baseline.
    \item \textbf{Group B (Level 2):} 4 task-oriented medium tools (search, fetch-and-summarize, write the review, append the bibliography) that chunk logical steps together.
    \item \textbf{Group C (Level 3):} A single monolithic macro-tool executing the entire literature review pipeline autonomously.
\end{itemize}

We hypothesize that Level 2 tools represent the optimal design compromise. By mitigating the cascading planning failures associated with micro-tools while preserving the adaptive reasoning lost in macro-tools, medium-abstraction tools should be at least as reliable as micro-tools at a lower cost, while still satisfying new requirements that a fixed macro-tool pipeline cannot (Section~\ref{sec:hypotheses}). 

To rigorously test this hypothesis, we benchmark five locally served models---three families of similar size and three sizes of one family---across 10 evaluation scenarios, measuring token consumption, model and tool calls, and execution time \cite{patil2025bfcl}. Furthermore, agent success is measured through ten deterministic binary checks that assess output structural integrity, APA citation formatting, the factual verification of arXiv identifiers, and evidence that every cited paper was actually read \cite{liu2024agentbench, yao2024taubench}. By mapping the trade-offs between tool granularity and model scale, this work provides actionable design guidelines for building robust autonomous agents on local infrastructure.